\FigureLayoutDeclare{img/2010/jiangxi_liberal/q21_fig1.png}{9.0cm}{47bp 71bp 0bp 17bp}

\chapter{2010年江西卷（文）}

\section{单选题}

\begin{problem}
对于实数 \(a\)，\(b\)，\(c\)，“\(a > b\)”是“\(ac^2 > bc^2\)”的（\quad）

\choices
    {充分不必要条件}
    {必要不充分条件}
    {充要条件}
    {既不充分也不必要条件}
\end{problem}

\begin{answer}
B
\end{answer}

\begin{problem}
若集合 \(A = \{x \mid |x| \leqslant 1\}\)，\(B = \{x \mid x \geqslant 0\}\)，则 \(A \cap B =\)（\quad）

\choices
    {\(\{x \mid -1 \leqslant x \leqslant 1\}\)}
    {\(\{x \mid x \geqslant 0\}\)}
    {\(\{x \mid 0 \leqslant x \leqslant 1\}\)}
    {\(\varnothing\)}
\end{problem}

\begin{answer}
C
\end{answer}

\begin{problem}
\((1 - x)^{10}\) 展开式中 \(x^{3}\) 项的系数为（\quad）

\choices
    {\(-720\)}
    {\(720\)}
    {\(120\)}
    {\(-120\)}
\end{problem}

\begin{answer}
D
\end{answer}

\begin{problem}
若函数 \(f(x) = ax^4 + bx^2 + c\) 满足 \(f'(1) = 2\)，则 \(f'(-1) =\)（\quad）

\choices
    {\(-1\)}
    {\(-2\)}
    {\(2\)}
    {\(0\)}
\end{problem}

\begin{answer}
B
\end{answer}

\begin{problem}
不等式 \(|x - 2| > x - 2\) 的解集是（\quad）

\choices
    {\((-\infty, 2)\)}
    {\((-\infty, +\infty)\)}
    {\((2, +\infty)\)}
    {\((-\infty, 2) \cup (2, +\infty)\)}
\end{problem}

\begin{answer}
A
\end{answer}

\begin{problem}
函数 \(y = \sin^{2}x + \sin x - 1\) 的值域为（\quad）

\choices
    {\([-1, 1]\)}
    {\(\left[-\frac{5}{4}, -1\right]\)}
    {\(\left[-\frac{5}{4}, 1\right]\)}
    {\(\left[-1, \frac{5}{4}\right]\)}
\end{problem}

\begin{answer}
C
\end{answer}

\begin{problem}
等比数列 \(\{a_{n}\}\) 中，\(|a_{1}| = 1\)，\(a_{5} = -8a_{2}\)，\(a_{5} > a_{2}\)，则 \(a_{n} =\)（\quad）

\choices
    {\((-2)^{n - 1}\)}
    {\(-(-2)^{n - 1}\)}
    {\((-2)^{n}\)}
    {\(-(-2)^{n}\)}
\end{problem}

\begin{answer}
A
\end{answer}

\begin{problem}
若函数 \(y = \frac{ax}{1 + x}\) 的图象关于直线 \(y = x\) 对称，则 \(a\) 为（\quad）

\choices
    {\(1\)}
    {\(-1\)}
    {\(\pm 1\)}
    {\(\text{任意实数}\)}
\end{problem}

\begin{answer}
B
\end{answer}

\begin{problem}
有 \(n\) 位同学参加某项选拔测试，每位同学能通过测试的概率都是 \(p\) （\(0 < p < 1\)），假设每位同学能否通过测试是相互独立的，则至少有一位同学通过测试的概率为（\quad）

\choices
    {\((1 - p)^n\)}
    {\(1 - p^n\)}
    {\(p^n\)}
    {\(1 - (1 - p)^n\)}
\end{problem}

\begin{answer}
D
\end{answer}

\begin{problem}
直线 \(y = kx + 3\) 与圆 \((x - 2)^2 + (y - 3)^2 = 4\) 相交于 \(M\)，\(N\) 两点，若 \(|MN| \geqslant 2\sqrt{3}\)，则 \(k\) 的取值范围是（\quad）

\choices
    {\(\left[-\frac{3}{4}, 0\right]\)}
    {\(\left[-\frac{\sqrt{3}}{3}, \frac{\sqrt{3}}{3}\right]\)}
    {\(\left[-\sqrt{3}, \sqrt{3}\right]\)}
    {\(\left[-\frac{2}{3}, 0\right]\)}
\end{problem}

\begin{answer}
B
\end{answer}

\begin{problem}
如图，\(M\) 是正方体 \(ABCD-A_{1}B_{1}C_{1}D_{1}\) 的棱 \(DD_{1}\) 的中点，给出下列命题：

\bitmapfigure[max width=0.35\linewidth]{img/2010/jiangxi_liberal/q11_fig1.png}

\begin{circlelist}
\item 过 \(M\) 点有且只有一条直线与直线 \(AB\)，\(B_{1}C_{1}\) 都相交；
\item 过 \(M\) 点有且只有一条直线与直线 \(AB\)，\(B_{1}C_{1}\) 都垂直；
\item 过 \(M\) 点有且只有一个平面与直线 \(AB\)，\(B_{1}C_{1}\) 都相交；
\item 过 \(M\) 点有且只有一个平面与直线 \(AB\)，\(B_{1}C_{1}\) 都平行.
\end{circlelist}
其中真命题是（\quad）

\choices
    {\circled{2}\circled{3}\circled{4}}
    {\circled{1}\circled{3}\circled{4}}
    {\circled{1}\circled{2}\circled{4}}
    {\circled{1}\circled{2}\circled{3}}
\end{problem}

\begin{answer}
C
\end{answer}

\begin{problem}
四位同学在同一个坐标系中分别选定了一个适当的区间，各自作出三个函数 \(y = \sin 2x\)，\(y = \sin \left(x + \frac{\pi}{6}\right)\)，\(y = \sin \left(x - \frac{\pi}{3}\right)\) 的图象如下. 结果发现其中有一位同学作出的图象有错误，那么有错误的图象是（\quad）

\choices
    {\choicebitmap[width=0.15\paperwidth]{img/2010/jiangxi_liberal/q11_optA_nolabel.png}}
    {\choicebitmap[width=0.15\paperwidth]{img/2010/jiangxi_liberal/q11_optB_nolabel.png}}
    {\choicebitmap[width=0.15\paperwidth]{img/2010/jiangxi_liberal/q11_optC_nolabel.png}}
    {\choicebitmap[width=0.15\paperwidth]{img/2010/jiangxi_liberal/q11_optD_nolabel.png}}
\end{problem}

\begin{answer}
C
\end{answer}

\section{填空题}

\begin{problem}
已知向量 \(\bs{a}\)，\(\bs{b}\) 满足 \(\lvert\bs{b}\rvert = 2\)，\(\bs{a}\) 与 \(\bs{b}\) 的夹角为 \(60^\circ\)，则 \(\bs{b}\) 在 \(\bs{a}\) 上的投影是\fillinblank{}.
\end{problem}

\begin{answer}
1
\end{answer}

\begin{problem}
将 \(5\) 位志愿者分成 \(3\) 组，其中两组各 \(2\) 人，另一组 \(1\) 人，分赴世博会的 \(3\) 个不同场馆服务，不同的分配方案有\fillinblank{}种.
\end{problem}

\begin{answer}
90
\end{answer}

\begin{problem}
点 \(A(x_{0},y_{0})\) 在双曲线 \(\frac{x^{2}}{4}-\frac{y^{2}}{32}=1\) 的右支上，若点 \(A\) 到右焦点的距离等于 \(2x_{0}\)，则 \(x_{0}=\)\fillinblank{}.
\end{problem}

\begin{answer}
2
\end{answer}

\begin{problem}
长方体 \(ABCD - A_{1}B_{1}C_{1}D_{1}\) 的顶点均在同一个球面上，\(AB = AA_{1} = 1\)，\(BC = \sqrt{2}\)，则 \(A\)，\(B\) 两点间的球面距离为\fillinblank{}.

\bitmapfigure[max width=0.42\linewidth]{img/2010/jiangxi_liberal/q16_fig1.png}
\end{problem}

\begin{answer}
\(\dfrac{\pi}{3}\)
\end{answer}

\section{解答题}

\begin{problem}
设函数 \(f(x) = 6x^{3} + 3(a + 2)x^{2} + 2ax\).
\begin{enumerate}
\item 若 \(f(x)\) 的两个极值点为 \(x_{1}\)，\(x_{2}\)，且 \(x_{1}x_{2} = 1\)，求实数 \(a\) 的值；
\item 是否存在实数 \(a\)，使得 \(f(x)\) 是 \((-\infty, +\infty)\) 上的单调函数？若存在，求出 \(a\) 的值；若不存在，说明理由.
\end{enumerate}
\end{problem}

\begin{solution}
\begin{enumerate}
\item 求导得 \(f'(x)=18x^{2}+6(a+2)x+2a\). 由于 \(f(x)\) 有两个极值点，所以这两个极值点的横坐标 \(x_{1}\)，\(x_{2}\) 是方程 \(f'(x)=0\) 的两个不同实根. 由韦达定理，
\[
x_{1}x_{2}=\frac{2a}{18}=\frac{a}{9}.
\]
由题设 \(x_{1}x_{2}=1\)，得 \(a=9\). 此时 \(f'(x)=6(3x^{2}+11x+3)\)，而二次式 \(3x^{2}+11x+3\) 的判别式为 \(11^{2}-4\times3\times3=85>0\)，故方程 \(f'(x)=0\) 确有两个不同实根，\(f(x)\) 确有两个极值点. 因此 \(a=9\).

\item 对任意实数 \(a\)，\(f'(x)\) 是二次项系数为正的二次函数，且
\[
\Delta=\bigl(6(a+2)\bigr)^{2}-4\times18\times2a=36(a^{2}+4)>0.
\]
所以方程 \(f'(x)=0\) 对任意实数 \(a\) 都有两个不同实根，设其为 \(r_{1}<r_{2}\). 由于二次项系数为正，\(f'(x)>0\) 在 \(r_{1}\)，\(r_{2}\) 的外侧成立，而 \(f'(x)<0\) 在 \((r_{1},r_{2})\) 内成立. 因此 \(f(x)\) 在 \((-\infty,r_{1})\) 上递增，在 \((r_{1},r_{2})\) 上递减，在 \((r_{2},+\infty)\) 上递增，不可能在 \((-\infty,+\infty)\) 上单调. 不存在满足条件的实数 \(a\).
\end{enumerate}
\end{solution}

\begin{problem}
某迷宫有三个通道，进入迷宫的每个人都要经过一个智能门. 首次到达此门，系统会随机（即等可能）为你打开一个通道. 若是 1 号通道，则需要 1 小时走出迷宫；若是 2 号，3 号通道，则分别需要 2 小时，3 小时返回智能门. 再次到达智能门时，系统会随机打开一个你未到过的通道，直至走出迷宫为止.
\begin{enumerate}
\item 求走出迷宫时恰好用了 \(1\) 小时的概率；
\item 求走出迷宫的时间超过 \(3\) 小时的概率.
\end{enumerate}
\end{problem}

\begin{solution}
\begin{enumerate}
\item 第一次到达智能门时，三个通道被打开的概率都为 \(\frac{1}{3}\). 只有第一次打开 \(1\) 号通道时能在恰好 \(1\) 小时后走出迷宫，所以 \(P(\text{恰好用时 }1\text{ 小时})=\frac{1}{3}\).

\item 设走出迷宫所用的时间为 \(T\). 若第一次打开 \(1\) 号通道，则 \(T=1\)，不满足 \(T>3\). 若第一次打开 \(2\) 号通道，经过 \(2\) 小时回到智能门，剩下的 \(1\) 号，\(3\) 号通道等可能被打开：打开 \(1\) 号通道时 \(T=2+1=3\)，打开 \(3\) 号通道后还要打开 \(1\) 号通道，此时 \(T=2+3+1=6\). 因此在第一次打开 \(2\) 号通道的条件下，\(P(T>3)=\frac{1}{2}\).

若第一次打开 \(3\) 号通道，经过 \(3\) 小时回到智能门；此后打开 \(1\) 号通道时 \(T=3+1=4\)，打开 \(2\) 号通道后还要打开 \(1\) 号通道时 \(T=3+2+1=6\)，两种情况都满足 \(T>3\). 因此
\[
P(T>3)=\frac{1}{3}\times\frac{1}{2}+\frac{1}{3}\times1=\frac{1}{2}.
\]
\end{enumerate}
\end{solution}

\begin{problem}
已知函数 \(f(x) = (1 + \cot x)\sin^{2}x - 2\sin \left(x + \frac{\pi}{4}\right)\sin \left(x - \frac{\pi}{4}\right)\).
\begin{enumerate}
\item 若 \(\tan \alpha = 2\)，求 \(f(\alpha)\)；
\item 若 \(x \in \left[\frac{\pi}{12}, \frac{\pi}{2}\right]\)，求 \(f(x)\) 的取值范围.
\end{enumerate}
\end{problem}

\begin{solution}
先化简函数. 函数有意义时 \(\sin x\ne0\)，于是 \((1+\cot x)\sin^{2}x=\sin^{2}x+\sin x\cos x\). 又由积化和差公式，
\[
2\sin\left(x+\frac{\pi}{4}\right)\sin\left(x-\frac{\pi}{4}\right)=-\cos2x.
\]
因此
\[
f(x)=\sin^{2}x+\sin x\cos x+\cos2x=\cos^{2}x+\sin x\cos x=\cos x(\sin x+\cos x).
\]

\begin{enumerate}
\item 由 \(\tan\alpha=2\)，得
\[
\cos^{2}\alpha=\frac{1}{1+\tan^{2}\alpha}=\frac{1}{5}.
\]
由 \(f(\alpha)=\cos^{2}\alpha(1+\tan\alpha)\)，得 \(f(\alpha)=\frac{1}{5}(1+2)=\frac{3}{5}\).

\item 在区间 \(\left[\frac{\pi}{12},\frac{\pi}{2}\right]\) 上，\(f(x)\) 有意义且连续. 求导得 \(f'(x)=\cos2x-\sin2x\). 因为 \(2x\in\left[\frac{\pi}{6},\pi\right]\)，方程 \(f'(x)=0\) 等价于 \(\tan2x=1\)，在该区间内的唯一解为 \(x=\frac{\pi}{8}\). 当 \(x\in\left[\frac{\pi}{12},\frac{\pi}{8}\right]\) 时，\(f'(x)\geqslant0\)；当 \(x\in\left[\frac{\pi}{8},\frac{\pi}{2}\right]\) 时，\(f'(x)\leqslant0\). 所以 \(f(x)\) 在前一段递增，在后一段递减，最大值在 \(x=\frac{\pi}{8}\) 处取得，且
\[
f\left(\frac{\pi}{8}\right)=\frac{1+\cos\frac{\pi}{4}}{2}+\frac{\sin\frac{\pi}{4}}{2}=\frac{1+\sqrt{2}}{2}.
\]
在给定区间内，\(\cos x\geqslant0\)，\(\sin x+\cos x>0\)，所以 \(f(x)\geqslant0\). 又 \(f\left(\frac{\pi}{2}\right)=0\)，故最小值为 \(0\). 因此 \(f(x)\) 的取值范围为 \(\left[0,\frac{1+\sqrt{2}}{2}\right]\).
\end{enumerate}
\end{solution}

\begin{problem}
如图，\(\triangle BCD\) 与 \(\triangle MCD\) 都是边长为 \(2\) 的正三角形，平面 \(MCD \perp\) 平面 \(BCD\)，\(AB \perp\) 平面 \(BCD\)，\(AB = 2\sqrt{3}\).
\begin{enumerate}
\item 求直线 \(AM\) 与平面 \(BCD\) 所成的角的大小；
\item 求平面 \(ACM\) 与平面 \(BCD\) 所成二面角的正弦值.
\end{enumerate}

\bitmapfigure[max width=0.42\linewidth]{img/2010/jiangxi_liberal/q20_fig1.png}
\end{problem}

\begin{solution}
建立空间直角坐标系，使平面 \(BCD\) 为 \(z=0\)，取 \(B=(0,0,0)\)，\(C=(2,0,0)\)，\(D=(1,\sqrt{3},0)\). 由 \(AB\perp\) 平面 \(BCD\) 且 \(AB=2\sqrt{3}\)，可取 \(A=(0,0,2\sqrt{3})\). 设 \(H\) 为 \(CD\) 的中点，则 \(H=\left(\frac{3}{2},\frac{\sqrt{3}}{2},0\right)\). 由于平面 \(MCD\) 垂直于平面 \(BCD\)，且 \(\triangle MCD\) 是边长为 \(2\) 的正三角形，故 \(MH\perp CD\)，\(MH=\sqrt{3}\)，从而可取 \(M=\left(\frac{3}{2},\frac{\sqrt{3}}{2},\sqrt{3}\right)\).

\begin{enumerate}
\item \(\overrightarrow{AM}=\left(\frac{3}{2},\frac{\sqrt{3}}{2},-\sqrt{3}\right)\)，所以
\[
|AM|=\sqrt{\frac{9}{4}+\frac{3}{4}+3}=\sqrt{6}.
\]
直线 \(AM\) 在平面 \(BCD\) 上的射影长度为
\[
\sqrt{\left(\frac{3}{2}\right)^{2}+\left(\frac{\sqrt{3}}{2}\right)^{2}}=\sqrt{3}
\]
直线 \(AM\) 垂直于平面 \(BCD\) 的分量长度也是 \(\sqrt{3}\). 设直线 \(AM\) 与平面 \(BCD\) 所成的角为 \(\theta\)，则
\[
\sin\theta=\frac{\sqrt{3}}{\sqrt{6}}=\frac{\sqrt{2}}{2}
\]
所以 \(\theta=45^\circ\).

\item 取平面 \(ACM\) 内的两个相交向量 \(\overrightarrow{AC}=(2,0,-2\sqrt{3})\)，\(\overrightarrow{AM}=\left(\frac{3}{2},\frac{\sqrt{3}}{2},-\sqrt{3}\right)\). 它们的一个法向量为
\[
\bs{n}=\overrightarrow{AC}\times\overrightarrow{AM}=(3,-\sqrt{3},\sqrt{3})
\]
平面 \(BCD\) 的一个法向量为 \(\bs{k}=(0,0,1)\)，且
\[
|\bs{n}|=\sqrt{15},\qquad
\frac{|\bs{n}\cdot\bs{k}|}{|\bs{n}||\bs{k}|}=\frac{\sqrt{3}}{\sqrt{15}}=\frac{\sqrt{5}}{5}
\]
设两平面所成的锐二面角为 \(\varphi\)，则 \(\cos\varphi=\frac{\sqrt{5}}{5}\)，所以
\[
\sin\varphi=\sqrt{1-\frac{1}{5}}=\frac{2\sqrt{5}}{5}.
\]
\end{enumerate}
\end{solution}

\begin{problem}
如图，已知抛物线 \(C_{1}: x^{2} + by = b^{2}\) 经过椭圆 \(C_{2}: \frac{x^{2}}{a^{2}} + \frac{y^{2}}{b^{2}} = 1\) （\(a > b > 0\)） 的两个焦点.
\begin{enumerate}
\item 求椭圆 \(C_2\) 的离心率；
\item 设 \(Q(3, b)\)，又 \(M\)，\(N\) 为 \(C_1\) 与 \(C_2\) 不在 \(y\) 轴上的两个交点，若 \(\triangle QMN\) 的重心在抛物线 \(C_1\) 上，求 \(C_1\) 和 \(C_2\) 的方程.
\end{enumerate}

\bitmapfigure[max width=0.42\linewidth]{img/2010/jiangxi_liberal/q21_fig1.png}
\end{problem}

\begin{solution}
设椭圆 \(C_{2}\) 的半焦距为 \(c\)，则其两个焦点为 \((\pm c,0)\)，且
\[
c^{2}=a^{2}-b^{2}.
\]
由于抛物线 \(C_{1}:x^{2}+by=b^{2}\) 经过这两个焦点，代入 \((\pm c,0)\) 得 \(c^{2}=b^{2}\). 因为 \(b>0\)，所以 \(c=b\). 结合 \(a>b>0\)，可得
\[
a^{2}=b^{2}+c^{2}=2b^{2}
\]
故椭圆 \(C_{2}\) 的离心率为
\[
e=\frac{c}{a}=\frac{b}{\sqrt{2}b}=\frac{\sqrt{2}}{2}.
\]

设 \(C_{1}\) 与 \(C_{2}\) 不在 \(y\) 轴上的两个交点为 \(M(-u,v)\)，\(N(u,v)\)，其中 \(u>0\). 由于两曲线的方程中只含 \(x^{2}\)，两交点的纵坐标相同. 将 \(N(u,v)\) 代入两曲线方程，得
\[
\begin{cases}
u^{2}+bv=b^{2},\\
\dfrac{u^{2}}{2}+v^{2}=b^{2}.
\end{cases}
\]
由第一式得 \(u^{2}=b^{2}-bv\)，代入第二式，得
\[
\frac{b^{2}-bv}{2}+v^{2}=b^{2}
\quad\Longleftrightarrow\quad
2v^{2}-bv-b^{2}=0
\quad\Longleftrightarrow\quad
(2v+b)(v-b)=0.
\]
若 \(v=b\)，由 \(u^{2}+bv=b^{2}\) 得 \(u=0\)，这与交点不在 \(y\) 轴上矛盾. 因此 \(v=-\frac{b}{2}\)，并且 \(u^{2}=b^{2}-b\left(-\frac{b}{2}\right)=\frac{3b^{2}}{2}\).

由 \(Q=(3,b)\)，三角形 \(QMN\) 的重心为
\[
G=\left(\frac{3-u+u}{3},\frac{b-\frac{b}{2}-\frac{b}{2}}{3}\right)=(1,0).
\]
重心 \(G\) 在抛物线 \(C_{1}\) 上，所以
\[
1+b\times0=b^{2}.
\]
因 \(b>0\)，得 \(b=1\). 于是
\(C_{1}:x^{2}+y=1\)，\(C_{2}:\frac{x^{2}}{2}+y^{2}=1\).
\end{solution}

\begin{problem}
正实数数列 \(\{a_{n}\}\) 中，\(a_{1} = 1\)，\(a_{2} = 5\)，且 \(\{a_{n}^{2}\}\) 成等差数列.
\begin{enumerate}
\item 证明：数列 \(\{a_{n}\}\) 中有无穷多项为无理数；
\item 当 \(n\) 为何值时，\(a_{n}\) 为整数，并求出使 \(a_{n} < 200\) 的所有整数项的和.
\end{enumerate}
\end{problem}

\begin{solution}
由 \(a_{1}=1\)，\(a_{2}=5\)，得
\[
a_{1}^{2}=1,\qquad a_{2}^{2}=25
\]
因为 \(\{a_{n}^{2}\}\) 是等差数列，所以公差为 \(25-1=24\)，从而
\[
a_{n}^{2}=1+24(n-1)=24n-23.
\]
又因为各项为正数，故
\[
a_{n}=\sqrt{24n-23}.
\]

\begin{enumerate}
\item 若 \(a_{n}\) 为有理数，设 \(a_{n}=\frac{p}{q}\)，其中 \(p\)，\(q\) 为互质的正整数，则
\[
p^{2}=(24n-23)q^{2}.
\]
若素数 \(r\) 整除 \(q\)，则由上式可知 \(r\) 整除 \(p^{2}\)，从而 \(r\) 整除 \(p\)，这与 \(p\)，\(q\) 互质矛盾. 因此 \(q=1\)，所以 \(a_{n}\) 为整数. 也就是说，\(a_{n}\) 不是整数时必为无理数.

假设数列中只有有限多项为无理数，则存在正整数 \(N\)，使得当 \(n\geqslant N\) 时，\(a_{n}=m_{n}\) 都是正整数. 因为 \(a_{n}=\sqrt{24n-23}\) 随 \(n\) 增大而增大，所以 \(m_{n+1}>m_{n}\). 于是
\[
(m_{n+1}-m_{n})(m_{n+1}+m_{n})=m_{n+1}^{2}-m_{n}^{2}=24.
\]
但 \(m_{n}\) 随 \(n\) 增大而无限增大，取充分大的 \(n\) 可使 \(m_{n+1}+m_{n}>24\). 此时 \(m_{n+1}-m_{n}\) 是正整数，故上式左边大于 \(24\)，矛盾. 因此数列 \(\{a_{n}\}\) 中有无穷多项为无理数.

\item 设 \(a_{n}=m\) 为正整数，则
\[
m^{2}=24n-23
\quad\Longleftrightarrow\quad
m^{2}\equiv1\pmod{24}.
\]
若 \(m\) 被 \(2\) 或 \(3\) 整除，则 \(m^{2}\not\equiv1\pmod{24}\). 反之，若 \(m\) 不被 \(2\)，\(3\) 整除，则 \(m\equiv1\) 或 \(5\pmod{6}\)，即 \(m=6k+1\) 或 \(m=6k-1\)，而这两类数的平方都模 \(24\) 同余于 \(1\). 因此正整数项对应于
\[
m=6k+1\quad(k\geqslant0),\qquad\text{或}\qquad m=6k-1\quad(k\geqslant1)
\]
相应的 \(n\) 分别为
\[
n=\frac{m^{2}+23}{24}=\frac{3k^{2}+k+2}{2}\quad(k\geqslant0)
\]
或
\[
n=\frac{m^{2}+23}{24}=\frac{3k^{2}-k+2}{2}\quad(k\geqslant1)
\]
这就是 \(a_{n}\) 为整数时 \(n\) 的全部取值.

当 \(a_{n}<200\) 时，对应的正整数 \(m<200\). 对于 \(m=6k+1\)，有 \(0\leqslant k\leqslant33\)；对于 \(m=6k-1\)，有 \(1\leqslant k\leqslant33\). 因此所有整数项的和为
\[
\sum_{k=0}^{33}(6k+1)+\sum_{k=1}^{33}(6k-1)
=\bigl(6\times\frac{33\times34}{2}+34\bigr)+\bigl(6\times\frac{33\times34}{2}-33\bigr)
=3400+3333=6733.
\]
\end{enumerate}
\end{solution}
